\chapter{电子有效理论的逐分量复核与自旋修正}
\label{app:electron-checks}

正文只保留得到 PINEM 主方程所需的最短严格链。本附录集中保存三类不会改变主线、但适合逐式核对的计算：平方 Dirac 算符中的自旋--场项、固定自旋一维双带显式消元、非相对论极限以及 Foldy--Wouthuysen 高阶修正。

\section{平方 Dirac 算符：自旋--场项从非对易性出现}

把 Dirac 方程平方可以看清“标量相对论色散”和“自旋--场耦合”分别来自哪里。定义
\begin{equation}
 \sigma^{\mu\nu}\equiv\frac{\ii}{2}[\gamma^\mu,\gamma^\nu],
 \qquad
 F_{\mu\nu}\equiv\partial_\mu A_\nu-\partial_\nu A_\mu.
 \label{eq:sigma-F-def}
\end{equation}
由式~\eqref{eq:covariant-pi} 直接计算
\begin{align}
 [\pi_\mu,\pi_\nu]
 &=[\ii\hbar\partial_\mu-q_{\rm e}A_\mu,
 \ii\hbar\partial_\nu-q_{\rm e}A_\nu]\\
 &=-\ii q_{\rm e}\hbar
 \left(\partial_\mu A_\nu-\partial_\nu A_\mu\right)\\
 &=-\ii q_{\rm e}\hbar F_{\mu\nu}.
 \label{eq:pi-commutator}
\end{align}
左乘共轭的一阶算符 $\gamma^\nu\pi_\nu+mc$，有
\begin{align}
0
&=(\gamma^\nu\pi_\nu+mc)(\gamma^\mu\pi_\mu-mc)\Psi\\
&=\left(\gamma^\nu\gamma^\mu\pi_\nu\pi_\mu-m^2c^2\right)\Psi.
\end{align}
利用
\begin{equation}
 \gamma^\nu\gamma^\mu
 =g^{\nu\mu}\mathbf 1_4-\ii\sigma^{\nu\mu}
\end{equation}
且 $\sigma^{\nu\mu}$ 反对称，只有 $\pi_\nu\pi_\mu$ 的反对称部分有贡献：
\begin{align}
\gamma^\nu\gamma^\mu\pi_\nu\pi_\mu
&=\pi^\mu\pi_\mu
-\frac{\ii}{2}\sigma^{\nu\mu}[\pi_\nu,\pi_\mu]\\
&=\pi^\mu\pi_\mu
-\frac{q_{\rm e}\hbar}{2}\sigma^{\mu\nu}F_{\mu\nu}.
\end{align}
因此得到精确的二阶 Dirac 方程
\begin{equation}
 \boxed{
 \left[
 \pi^\mu\pi_\mu-m^2c^2
 -\frac{q_{\rm e}\hbar}{2}\sigma^{\mu\nu}F_{\mu\nu}
 \right]\Psi=0.}
 \label{eq:squared-dirac}
\end{equation}
这里第一项给出 Klein--Gordon 型相对论运动学，最后一项则是纯粹由协变动量不对易产生的自旋场项。用
\begin{equation}
 \bm\Sigma=
 \begin{pmatrix}
 \bm\sigma&0\\0&\bm\sigma
 \end{pmatrix},
 \qquad
 \sigma^{0i}=\ii\alpha_i,
 \qquad
 \sigma^{ij}=\epsilon_{ijk}\Sigma_k,
\end{equation}
以及本约定下
\begin{equation}
 F_{0i}=\frac{E_i}{c},
 \qquad
 F_{ij}=-\epsilon_{ijk}B_k,
\end{equation}
可把收缩显式写成
\begin{equation}
 \sigma^{\mu\nu}F_{\mu\nu}
 =\frac{2\ii}{c}\bm\alpha\cdot\bm E
 -2\bm\Sigma\cdot\bm B.
 \label{eq:sigmaF-EB}
\end{equation}
所以式~\eqref{eq:squared-dirac} 亦可写成
\begin{equation}
 \left[
 \pi^\mu\pi_\mu-m^2c^2
 +q_{\rm e}\hbar\bm\Sigma\cdot\bm B
 -\ii\frac{q_{\rm e}\hbar}{c}\bm\alpha\cdot\bm E
 \right]\Psi=0.
 \label{eq:squared-dirac-EB}
\end{equation}
这一步给出正能投影中磁矩、自旋--轨道与场梯度修正的母结构，也说明若只讨论纵向电场主导且固定自旋的 PINEM，相应自旋项为何可以作为受控高阶修正而不是凭经验删除。
式~\eqref{eq:squared-dirac-EB} 已经说明：若后文忽略显式自旋翻转或场梯度修正，必须把它视为从这个精确母式受控删项，而不能把自由电子一开始就当成标量粒子。现在才把外场关闭，建立后续正能投影所需要的自由本征基底。


\section{固定自旋的一维 Dirac 双带逐分量复核}
\label{app:dirac-two-band-check}

正文采用三维正能投影与质量张量建立包络方程，因为这条路线最短、规范结构也最清楚。本附录不建立新的物理模型，而是把同一结论在固定传播方向和固定自旋扇区中逐分量重算一次。目的只有两个：检查负能小分量的符号和快相位；逐项核对 $\partial_z^2f$、$A_z\partial_zf$、$(\partial_zA_z)f$ 与 $A_z^2f$ 的系数。读者若只关心 PINEM 主线，可以跳过本附录而不影响正文闭合。

\subsection{纵向约化：从四分量 Dirac 到固定自旋的两分量模型}
PINEM 中选中心电子动量
\begin{equation}
 \bm p_0=p_0\hat{\bm z},
 \qquad
 p_0=\hbar k_0=\gamma_0mv_0,
 \qquad
 \mathcal E_0=\gamma_0mc^2=\hbar\Omega_0.
\end{equation}
若横向电子动量宽度和光场引起的横向机械动量都远小于纵向载流动量，
\begin{equation}
 \boxed{
 \frac{\Delta p_\perp}{p_0}\ll1,
 \qquad
 \frac{|q_{\rm e}|A_\perp}{p_0}\ll1,
 \qquad
 \frac{\hbar|\bm\nabla_\perp f|}{p_0|f|}\ll1,}
 \label{eq:longitudinal-reduction-condition}
\end{equation}
则主动力学可投影到 $z$ 方向。再取 $\sigma_z\xi_s=s\xi_s$，$s=\pm1$。对 $\bm p=p\hat{\bm z}$，定义两个无量纲系数
\begin{equation}
 C_\gamma\equiv\sqrt{\frac{\gamma+1}{2\gamma}},
 \qquad
 S_\gamma\equiv\sqrt{\frac{\gamma-1}{2\gamma}},
 \qquad
 C_\gamma^2+S_\gamma^2=1.
 \label{eq:CG-SG}
\end{equation}
于是正能态与同一动量处的负能本征态分别为
\begin{equation}
 u_s(p\hat{\bm z})=
 \begin{pmatrix}
 C_\gamma\xi_s\\ sS_\gamma\xi_s
 \end{pmatrix},
 \qquad
 w_s(p\hat{\bm z})=
 \begin{pmatrix}
 -sS_\gamma\xi_s\\ C_\gamma\xi_s
 \end{pmatrix}.
 \label{eq:longitudinal-fourspinors}
\end{equation}
纵向势 $A_z$ 不混合 $\sigma_z$ 的两个自旋扇区，因此固定任意一个 $s$ 后都得到同构的二维问题。把固定 Pauli 自旋因子 $\xi_s$ 提出，并允许对负能基底作一个整体相位选择，可写成
\begin{equation}
 \boxed{
 \up(k)=
 \begin{pmatrix}
 C_\gamma\\ S_\gamma
 \end{pmatrix},
 \qquad
 \um(k)=
 \begin{pmatrix}
 -S_\gamma\\ C_\gamma
 \end{pmatrix}.}
 \label{eq:uplus}
\end{equation}
这里 $p=\hbar k$，$\gamma=\sqrt{1+\hbar^2k^2/(m^2c^2)}$。在这个固定自旋扇区中，$\alpha_z$ 的作用等价于二维 $\sigma_x$。逐项计算
\begin{align}
 \up{}^\dagger\sigma_x\up
 &=2C_\gamma S_\gamma
 =\frac{\sqrt{\gamma^2-1}}{\gamma}
 =\frac{v}{c},\\
 \um{}^\dagger\sigma_x\um
 &=-2C_\gamma S_\gamma=-\frac{v}{c},\\
 \up{}^\dagger\sigma_x\um
 &=C_\gamma^2-S_\gamma^2=\frac1\gamma,\\
 \um{}^\dagger\sigma_x\up&=\frac1\gamma.
 \label{eq:sigmax-matrix-elements}
\end{align}

在同一纵向近似下取 $\phi=0$、$\bm A\simeq A_z(z,t)\hat{\bm z}$。固定一个自旋扇区后，四分量 Dirac 方程严格约化为二维方程
\begin{equation}
 \boxed{
 \ii\hbar\frac{\partial\Psi}{\partial t}
 =\left[c\sigma_x\left(-\ii\hbar\frac{\partial}{\partial z}-q_{\rm e}A_z\right)
 +\sigma_zmc^2\right]\Psi.}
 \label{eq:dirac}
\end{equation}
这里的 $\Psi$ 已是固定 $\sigma_z$ 自旋扇区中的二分量振幅；若式~\eqref{eq:longitudinal-reduction-condition} 不成立，就必须回到四分量三维式~\eqref{eq:dirac-3d}，不能继续使用这一约化。

下面还需要中心动量附近自旋量随 $k$ 的变化。先由
\[
 \gamma(k)=\sqrt{1+\frac{\hbar^2k^2}{m^2c^2}}
\]
逐项求导：
\begin{align}
 \frac{\dd\gamma}{\dd k}
 &=\frac{\hbar^2k}{\gamma m^2c^2}
 =\frac{\hbar p}{\gamma m^2c^2}
 =\frac{\hbar v}{mc^2}.
 \label{eq:dgammak}
\end{align}
由
\[
 C_\gamma^2=\frac12\left(1+\frac1\gamma\right),
 \qquad
 S_\gamma^2=\frac12\left(1-\frac1\gamma\right)
\]
得
\begin{align}
 2C_\gamma\frac{\dd C_\gamma}{\dd k}
 &=-\frac{1}{2\gamma^2}\frac{\dd\gamma}{\dd k},\\
 2S_\gamma\frac{\dd S_\gamma}{\dd k}
 &=+\frac{1}{2\gamma^2}\frac{\dd\gamma}{\dd k}.
\end{align}
再用 $v/c=2C_\gamma S_\gamma$，可化为
\begin{align}
 \frac{\dd C_\gamma}{\dd k}
 &=-\frac{\hbar}{2\gamma^2mc}S_\gamma,\\
 \frac{\dd S_\gamma}{\dd k}
 &=+\frac{\hbar}{2\gamma^2mc}C_\gamma.
\end{align}
故
\begin{equation}
 \boxed{
 \frac{\partial\up(k)}{\partial k}
 =\frac{\hbar}{2\gamma^2mc}\um(k),
 \qquad
 \frac{\partial\um(k)}{\partial k}
 =-\frac{\hbar}{2\gamma^2mc}\up(k).}
 \label{eq:spinor-derivatives}
\end{equation}

\subsection{固定中心基底中的窄带正能波包}

先不指定包络形状。设纵向正能波包的谱振幅为 $a(k)$，归一化为
\begin{equation}
\int_{-\infty}^{+\infty}|a(k)|^2\,\dd k=1,
\end{equation}
且主要支撑在 $|k-k_0|\lesssim\Delta k$。初态完全由正能本征态组成：
\begin{equation}
\Psi(z,0)
=\frac{1}{\sqrt{2\pi}}
\int_{-\infty}^{+\infty}\!\dd k\,
 a(k)\up(k)\ee^{\ii kz}.
\label{eq:general-positive-packet}
\end{equation}
窄带条件不只是 $\Delta k\ll k_0$，还应直接控制自旋量 Taylor 余项：
\begin{equation}
\boxed{
\Delta k\left\|\left.\partial_k\up\right|_{k_0}\right\|\ll1,
\qquad
\frac{(\Delta k)^2}{2}
\left\|\left.\partial_k^2\up\right|_{k_0}\right\|
\ll
\Delta k\left\|\left.\partial_k\up\right|_{k_0}\right\|.}
\label{eq:spinor-taylor-condition}
\end{equation}
于是
\begin{equation}
\up(k)
=\up_0+(k-k_0)
\left.\frac{\partial\up}{\partial k}\right|_{k_0}
+O[(k-k_0)^2],
\qquad
\up_0\equiv\up(k_0).
\end{equation}
定义去掉中心载波后的任意标量包络
\begin{equation}
f_0(z)
\equiv
\frac{1}{\sqrt{2\pi}}
\int_{-\infty}^{+\infty}\!\dd k\,
 a(k)\ee^{\ii(k-k_0)z}.
\label{eq:general-envelope-Fourier}
\end{equation}
由
\begin{equation}
\partial_zf_0(z)
=\frac{\ii}{\sqrt{2\pi}}
\int (k-k_0)a(k)\ee^{\ii(k-k_0)z}\dd k
\end{equation}
可得
\begin{equation}
\frac{1}{\sqrt{2\pi}}
\int (k-k_0)a(k)\ee^{\ii(k-k_0)z}\dd k
=-\ii\partial_zf_0(z).
\end{equation}
将其代入式~\eqref{eq:general-positive-packet}，再使用式~\eqref{eq:spinor-derivatives}，得到与包络形状无关的一阶结果
\begin{equation}
\boxed{
\Psi(z,0)
\simeq
\ee^{\ii k_0z}
\left[
 f_0(z)\up_0
-\ii\frac{\hbar}{2\gamma_0^2mc}
 (\partial_zf_0)\um_0
\right].}
\label{eq:general-fixed-basis-packet}
\end{equation}
这里出现的 $\um_0$ 不表示波包中真实产生了正电子。式~\eqref{eq:general-positive-packet} 在每个 $k$ 上都严格位于正能子空间；$\um_0$ 只是在把不同 $k$ 的正能自旋量统一展开到中心基底 $\{\up_0,\um_0\}$ 时，由 $\partial_k\up\propto\um$ 产生的基底变化项。高斯波包对应于把 $a(k)$ 进一步取为高斯函数。

\subsection{两能带包络方程与快速负能分量的显式消去}

在中心动量 $k_0$ 处，令
\[
 \up_0\equiv\up(k_0),
 \qquad
 \um_0\equiv\um(k_0),
\]
并写
\begin{equation}
 \Psi(z,t)=\ee^{\ii k_0z}
 \left[
 f^+(z,t)\up_0\ee^{-\ii\Omega_0t}
 +f^-(z,t)\um_0\ee^{+\ii\Omega_0t}
 \right].
 \label{eq:envelope-ansatz}
\end{equation}
其中 $+\mathcal E_0$ 分量的自由时间因子为 $\ee^{-\ii\Omega_0t}$，$-\mathcal E_0$ 分量的自由时间因子为 $\ee^{+\ii\Omega_0t}$。

对任意包络 $f(z,t)$，动量算符作用为
\begin{align}
 -\ii\hbar\partial_z\left[f\widehat u\ee^{\ii k_0z}\right]
 &=-\ii\hbar\left[(\partial_z f)\widehat u\ee^{\ii k_0z}
 +\ii k_0f\widehat u\ee^{\ii k_0z}\right]\\
 &=\left[p_0f-\ii\hbar\partial_z f\right]
 \widehat u\ee^{\ii k_0z}.
 \label{eq:momentum-envelope-action}
\end{align}
把式 \eqref{eq:envelope-ansatz} 代入式 \eqref{eq:dirac}。其中 $p_0$ 和 $mc^2$ 的自由部分分别由
\[
 (c\sigma_xp_0+\sigma_zmc^2)\up_0=+\mathcal E_0\up_0,
 \qquad
 (c\sigma_xp_0+\sigma_zmc^2)\um_0=-\mathcal E_0\um_0
\]
与快速相位 $\ee^{\mp\ii\Omega_0t}$ 抵消。剩余项满足
\begin{align}
 \ii\hbar\Bigl[
 (\partial_t f^+)\up_0\ee^{-\ii\Omega_0t}
 +(\partial_t f^-)\um_0\ee^{+\ii\Omega_0t}
 \Bigr]
 ={}&-\ii\hbar c\sigma_x\Bigl[
 (\partial_z f^+)\up_0\ee^{-\ii\Omega_0t}
 +(\partial_z f^-)\um_0\ee^{+\ii\Omega_0t}
 \Bigr]\notag\\
 &-q_{\rm e}cA_z\sigma_x\Bigl[
 f^+\up_0\ee^{-\ii\Omega_0t}
 +f^-\um_0\ee^{+\ii\Omega_0t}
 \Bigr].
 \label{eq:after-free-cancel}
\end{align}

左乘 $\widehat u_0^{+\dagger}\ee^{+\ii\Omega_0t}$，并用式 \eqref{eq:sigmax-matrix-elements} 在 $k_0$ 处的结果
\[
 \widehat u_0^{+\dagger}\sigma_x\up_0=\frac{v_0}{c},
 \qquad
 \widehat u_0^{+\dagger}\sigma_x\um_0=\frac1{\gamma_0},
\]
得到
\begin{align}
 \ii\hbar\partial_t f^+
 ={}&-\ii\hbar v_0\partial_z f^+
 -\ii\hbar\frac{c}{\gamma_0}\ee^{+2\ii\Omega_0t}\partial_z f^-\\
 &-q_{\rm e}v_0A_zf^+
 -q_{\rm e}\frac{c}{\gamma_0}A_z\ee^{+2\ii\Omega_0t}f^-.
\end{align}
两边除以 $\ii\hbar$，再把传播项移到左边：
\begin{equation}
 \boxed{
 \partial_t f^+
 +v_0\partial_z f^+
 +\frac{c}{\gamma_0}\ee^{+2\ii\Omega_0t}\partial_z f^-
 =\ii\frac{q_{\rm e}}{\hbar}A_z
 \left[v_0f^+
 +\frac{c}{\gamma_0}\ee^{+2\ii\Omega_0t}f^-\right].}
 \label{eq:coupled-plus}
\end{equation}

负能投影需要同样保留全部矩阵元和快速相位。左乘
$\widehat u_0^{-\dagger}\ee^{-\ii\Omega_0t}$。由正交归一性，式~\eqref{eq:after-free-cancel}
左端只剩
\[
 \widehat u_0^{-\dagger}\ee^{-\ii\Omega_0t}
 \Bigl[\ii\hbar(\partial_t f^-)\um_0\ee^{+\ii\Omega_0t}\Bigr]
 =\ii\hbar\partial_t f^-.
\]
空间导数项需要两个矩阵元
\[
 \widehat u_0^{-\dagger}\sigma_x\um_0=-\frac{v_0}{c},
 \qquad
 \widehat u_0^{-\dagger}\sigma_x\up_0=\frac1{\gamma_0}.
\]
因此正能分量投影到负能带时会留下快速相位
\[
 \widehat u_0^{-\dagger}\ee^{-\ii\Omega_0t}
 \sigma_x\up_0\ee^{-\ii\Omega_0t}
 =\frac1{\gamma_0}\ee^{-2\ii\Omega_0t},
\]
而负能分量自身没有额外快相位。于是传播部分逐项为
\begin{align}
 &-\ii\hbar c\,
 \widehat u_0^{-\dagger}\ee^{-\ii\Omega_0t}\sigma_x
 \Bigl[(\partial_z f^+)\up_0\ee^{-\ii\Omega_0t}
 +(\partial_z f^-)\um_0\ee^{+\ii\Omega_0t}\Bigr]\notag\\
 &\qquad=
 -\ii\hbar\frac{c}{\gamma_0}\ee^{-2\ii\Omega_0t}\partial_z f^+
 +\ii\hbar v_0\partial_z f^- .
\end{align}
相互作用项同样逐项给出
\begin{align}
 &-q_{\rm e}cA_z\,
 \widehat u_0^{-\dagger}\ee^{-\ii\Omega_0t}\sigma_x
 \Bigl[f^+\up_0\ee^{-\ii\Omega_0t}
 +f^-\um_0\ee^{+\ii\Omega_0t}\Bigr]\notag\\
 &\qquad=
 -q_{\rm e}\frac{c}{\gamma_0}A_z\ee^{-2\ii\Omega_0t}f^+
 +q_{\rm e}v_0A_zf^- .
\end{align}
把两部分相加，得到
\begin{align}
 \ii\hbar\partial_t f^-
 ={}&+\ii\hbar v_0\partial_z f^-
 -\ii\hbar\frac{c}{\gamma_0}\ee^{-2\ii\Omega_0t}\partial_z f^+\\
 &+q_{\rm e}v_0A_zf^-
 -q_{\rm e}\frac{c}{\gamma_0}A_z\ee^{-2\ii\Omega_0t}f^+.
\end{align}
再把 $\partial_z f^-$ 项移到左端并除以 $\ii\hbar$，得到
\begin{equation}
 \boxed{
 \partial_t f^-
 -v_0\partial_z f^-
 +\frac{c}{\gamma_0}\ee^{-2\ii\Omega_0t}\partial_z f^+
 =\ii\frac{q_{\rm e}}{\hbar}A_z
 \left[-v_0f^-
 +\frac{c}{\gamma_0}\ee^{-2\ii\Omega_0t}f^+\right].}
 \label{eq:coupled-minus}
\end{equation}

式 \eqref{eq:coupled-plus}--\eqref{eq:coupled-minus} 是在固定 $k_0$ 自旋量基底下的耦合包络方程。为了求随正能波包向 $+z$ 传播的那一小部分 $f^-$，令
\begin{equation}
 f^-(z,t)=h(z,t)\ee^{-2\ii\Omega_0t}.
 \label{eq:fminus-h}
\end{equation}
于是
\[
 \partial_t f^-=(\partial_t h-2\ii\Omega_0h)\ee^{-2\ii\Omega_0t},
 \qquad
 \partial_z f^-=(\partial_z h)\ee^{-2\ii\Omega_0t}.
\]
代入式 \eqref{eq:coupled-minus}，并约去公共因子 $\ee^{-2\ii\Omega_0t}$：
\begin{equation}
 \partial_t h-2\ii\Omega_0h-v_0\partial_z h
 +\frac{c}{\gamma_0}\partial_z f^+
 =\ii\frac{q_{\rm e}}{\hbar}A_z
 \left[-v_0h+\frac{c}{\gamma_0}f^+\right].
 \label{eq:h-exact}
\end{equation}

这里的高频近似不是把所有导数简单设为零，而是要求 $h$ 及驱动项在时间尺度 $1/(2\Omega_0)$ 上缓慢变化。逐项写成
\begin{equation}
 \boxed{
 \frac{|\partial_t h|}{2\Omega_0|h|}\ll1,
 \qquad
 \frac{v_0|\partial_z h|}{2\Omega_0|h|}\ll1,
 \qquad
 \frac{|q_{\rm e}|v_0|A_z|}{2\hbar\Omega_0}\ll1.}
 \label{eq:HFconditions-h}
\end{equation}
为保证由 $f^+$ 和 $A_z$ 构成的驱动本身也是慢变量，可用更直接的充分条件
\begin{equation}
 \boxed{
 \Omega_{\rm slow}
 \equiv
 \max\left(
 \left|\partial_t\ln f^+\right|,
 v_0\left|\partial_z\ln f^+\right|,
 \left|\partial_t\ln A_z\right|,
 v_0\left|\partial_z\ln A_z\right|,
 \frac{|q_{\rm e}|v_0|A_z|}{\hbar}
 \right)
 \ll2\Omega_0.}
 \label{eq:HFcondition-compact}
\end{equation}
在这些条件下，式 \eqref{eq:h-exact} 中与 $-2\ii\Omega_0h$ 相比的 $\partial_th$、$v_0\partial_zh$ 和 $A_zh$ 均为高阶小量。保留主平衡：
\begin{align}
 -2\ii\Omega_0h
 +\frac{c}{\gamma_0}\partial_z f^+
 &\simeq
 \ii\frac{q_{\rm e}}{\hbar}A_z\frac{c}{\gamma_0}f^+,\\
 -2\ii\Omega_0h
 &\simeq
 \frac{c}{\gamma_0}
 \left[-\partial_z f^+
 +\ii\frac{q_{\rm e}}{\hbar}A_zf^+\right].
\end{align}
故
\begin{equation}
 h\simeq
 \frac{1}{-2\ii\Omega_0}\frac{c}{\gamma_0}
 \left[-\partial_z f^+
 +\ii\frac{q_{\rm e}}{\hbar}A_zf^+\right].
\end{equation}
代回式 \eqref{eq:fminus-h}：
\begin{equation}
 \boxed{
 f^-(z,t)\simeq
 \frac{1}{-2\ii\Omega_0}\frac{c}{\gamma_0}
 \left[-\partial_z f^+
 +\ii\frac{q_{\rm e}}{\hbar}A_zf^+\right]
 \ee^{-2\ii\Omega_0t}.}
 \label{eq:fminus-approx}
\end{equation}
再用 $\hbar\Omega_0=\gamma_0mc^2$，得到更方便的形式
\begin{equation}
 f^-\simeq
 \frac{\ii\hbar}{2\gamma_0^2mc}
 \left[-\partial_z f^+
 +\ii\frac{q_{\rm e}}{\hbar}A_zf^+\right]
 \ee^{-2\ii\Omega_0t}.
 \label{eq:fminus-approx2}
\end{equation}

现在把式 \eqref{eq:fminus-approx2} 完整代回正能方程 \eqref{eq:coupled-plus}。为缩短中间式，记 $f\equiv f^+$。首先有
\begin{align}
 \frac{c}{\gamma_0}\ee^{2\ii\Omega_0t}f^-
 &=\frac{c}{\gamma_0}
 \frac{\ii\hbar}{2\gamma_0^2mc}
 \left[-f_z+\ii\frac{q_{\rm e}}{\hbar}A_zf\right]\\
 &=\frac{\ii\hbar}{2\gamma_0^3m}
 \left[-f_z+\ii\frac{q_{\rm e}}{\hbar}A_zf\right]\\
 &=-\frac{\ii\hbar}{2\gamma_0^3m}f_z
 -\frac{q_{\rm e}A_z}{2\gamma_0^3m}f.
 \label{eq:Bdef}
\end{align}
对 $z$ 求导：
\begin{align}
 \frac{c}{\gamma_0}\ee^{2\ii\Omega_0t}\partial_zf^-
 &=\partial_z\left(
 -\frac{\ii\hbar}{2\gamma_0^3m}f_z
 -\frac{q_{\rm e}A_z}{2\gamma_0^3m}f
 \right)\\
 &=-\frac{\ii\hbar}{2\gamma_0^3m}f_{zz}
 -\frac{q_{\rm e}}{2\gamma_0^3m}(\partial_zA_z)f
 -\frac{q_{\rm e}A_z}{2\gamma_0^3m}f_z.
 \label{eq:Bz}
\end{align}
正能方程右边的交叉项为
\begin{align}
 \ii\frac{q_{\rm e}}{\hbar}A_z
 \frac{c}{\gamma_0}\ee^{2\ii\Omega_0t}f^-
 & =\ii\frac{q_{\rm e}}{\hbar}A_z
 \left[-\frac{\ii\hbar}{2\gamma_0^3m}f_z
 -\frac{q_{\rm e}A_z}{2\gamma_0^3m}f\right]\\
 &=\frac{q_{\rm e}A_z}{2\gamma_0^3m}f_z
 -\ii\frac{q_{\rm e}^2A_z^2}{2\hbar\gamma_0^3m}f\\
 &=\frac{q_{\rm e}A_z}{2\gamma_0^3m}f_z
 +\frac{q_{\rm e}^2A_z^2}{2\ii\hbar\gamma_0^3m}f.
 \label{eq:crossrhs}
\end{align}
把式 \eqref{eq:Bz}、\eqref{eq:crossrhs} 代入式 \eqref{eq:coupled-plus}：
\begin{align}
 f_t+v_0f_z
 &-\frac{\ii\hbar}{2\gamma_0^3m}f_{zz}
 -\frac{q_{\rm e}}{2\gamma_0^3m}(\partial_zA_z)f
 -\frac{q_{\rm e}A_z}{2\gamma_0^3m}f_z\notag\\
 &\simeq
 \ii\frac{q_{\rm e}v_0}{\hbar}A_zf
 +\frac{q_{\rm e}A_z}{2\gamma_0^3m}f_z
 +\frac{q_{\rm e}^2A_z^2}{2\ii\hbar\gamma_0^3m}f.
\end{align}
将左边最后两项移到右边，两个 $A_zf_z$ 的半项相加，得到
\begin{equation}
 \boxed{
 \begin{aligned}
 \partial_t f^+
 +v_0\partial_z f^+
 -\frac{\ii\hbar}{2\gamma_0^3m}\partial_z^2f^+
 \simeq{}&
 \ii\frac{q_{\rm e}v_0}{\hbar}A_zf^+
 +\frac{q_{\rm e}A_z}{\gamma_0^3m}\partial_z f^+\\
 &+\frac{q_{\rm e}^2A_z^2}{2\ii\hbar\gamma_0^3m}f^+
 +\frac{q_{\rm e}}{2\gamma_0^3m}(\partial_zA_z)f^+.
 \end{aligned}}
 \label{eq:rel-master-explicit-check}
\end{equation}
这就是由 Dirac 方程得到的正能包络相对论主方程。它的纵向色散系数可写成
\[
 -\frac{\ii\hbar}{2m_\parallel}\partial_z^2f^+,
 \qquad
 \boxed{m_\parallel=\gamma_0^3m},
\]
即相对论纵向质量。

式~\eqref{eq:rel-master-explicit-check} 与正文由质量张量直接得到的式~\eqref{eq:rel-master} 完全相同。应当注意，$\partial_zA_z$ 项本身不是可单独赋予物理意义的 ``散度修正''，也不能在一般纳米近场里仅凭 $\bm\nabla\cdot\bm A=0$ 就把它删去。它来自 Hermitian 算符
\(
(-\ii\hbar\partial_z-q_{\rm e}A_z)^2
\)
的算符次序。真正规范协变的对象是未展开的式~\eqref{eq:rel-envelope-1d-covariant}，或三维的式~\eqref{eq:rel-envelope-tensor}。若选择 Coulomb gauge，条件为
\begin{equation}
\bm\nabla\cdot\bm A=0,
\label{eq:coulomb-gauge}
\end{equation}
但在含横向分量的近场中这并不等价于 $\partial_zA_z=0$。后面进入一阶 PINEM 极限时，$\partial_zA_z$ 与 $A_z\partial_z$、$A_z^2$、色散项一起作为二阶有效动力学整体舍去，而不是靠规范选择单独消失。


在附录的逐分量复核中，需要把正文已经得到的边带结构重新写回正、负能自旋量。为此只需使用一般 Fourier 分解
\begin{equation}
 g^+(z',+\infty)
 =\sum_{n=-\infty}^{+\infty}
 \ee^{\ii n\kappa z'}g_n^+(z'),
 \label{eq:g-sideband-sum}
\end{equation}
其中 $g_n^+(z')$ 是正文式~\eqref{eq:Pn-beta-standard}（连续波）或式~\eqref{eq:pulse-local-sideband-amplitude}（有限脉冲）对应的第 $n$ 阶包络振幅。

相互作用结束后 $A_z\to0$。式 \eqref{eq:fminus-approx} 变为
\begin{equation}
 f^-(z,t)
 \simeq
 -\ii\frac{\hbar}{2\gamma_0^2mc}
 \partial_zf^+(z,t)\ee^{-2\ii\Omega_0t}.
 \label{eq:fminus-free}
\end{equation}
将它代回式 \eqref{eq:envelope-ansatz}：
\begin{align}
 \Psi(z,t)
 &\simeq\ee^{\ii k_0z}
 \left[
 f^+\up_0\ee^{-\ii\Omega_0t}
 -\ii\frac{\hbar}{2\gamma_0^2mc}
 (\partial_zf^+)\ee^{-2\ii\Omega_0t}
 \um_0\ee^{+\ii\Omega_0t}
 \right]\\
 &=\left[
 f^+(z,t)\up_0
 -\ii\frac{\hbar}{2\gamma_0^2mc}\partial_zf^+(z,t)\um_0
 \right]
 \ee^{\ii(k_0z-\Omega_0t)}.
 \label{eq:psi-after-interaction}
\end{align}
由于 $f^+(z,t)=g^+(z-v_0t,+\infty)$，且 $z'=z-v_0t$，有 $\partial_z=\partial_{z'}$。代入式 \eqref{eq:g-sideband-sum}：
\begin{align}
 \partial_{z'}g^+(z',+\infty)
 &=\sum_n\partial_{z'}\left[
 \ee^{\ii n\kappa z'}g_n^+(z')\right]\\
 &=\sum_n\ee^{\ii n\kappa z'}
 \left[\ii n\kappa g_n^+(z')+\partial_{z'}g_n^+(z')\right].
 \label{eq:g-derivative-sidebands}
\end{align}
所以
\begin{align}
 \Psi(z,t)
 \simeq\sum_n\Bigg\{
 &g_n^+(z')\up_0
 -\ii\frac{\hbar}{2\gamma_0^2mc}
 \left[\ii n\kappa g_n^+(z')+\partial_{z'}g_n^+(z')\right]\um_0
 \Bigg\}\notag\\
 &\times\ee^{\ii n\kappa z'}\ee^{\ii(k_0z-\Omega_0t)}.
\end{align}
因为 $-\ii\times\ii=+1$，第一部分变成
\begin{align}
 \Psi(z,t)
 \simeq\sum_n\Bigg\{
 &g_n^+(z')
 \left[\up_0+n\kappa\frac{\hbar}{2\gamma_0^2mc}\um_0\right]
 -\ii\frac{\hbar}{2\gamma_0^2mc}
 (\partial_{z'}g_n^+)\um_0
 \Bigg\}\notag\\
 &\times\ee^{\ii n\kappa z'}\ee^{\ii(k_0z-\Omega_0t)}.
 \label{eq:psi-before-spinor-repack}
\end{align}

定义
\begin{equation}
 \delta k_n\equiv n\kappa=n\frac{\omega}{v_0},
 \qquad
 k_n\equiv k_0+\delta k_n.
 \label{eq:kn-def}
\end{equation}
由式 \eqref{eq:spinor-derivatives}，在 $k_0$ 处
\[
 \left.\frac{\partial\up}{\partial k}\right|_{k_0}
 =\frac{\hbar}{2\gamma_0^2mc}\um_0.
\]
因此
\begin{align}
 \up(k_n)
 &=\up(k_0)+\delta k_n
 \left.\frac{\partial\up}{\partial k}\right|_{k_0}
 +\frac{\delta k_n^2}{2}
 \left.\frac{\partial^2\up}{\partial k^2}\right|_{k_0}+\cdots\\
 &\simeq
 \up_0+n\kappa\frac{\hbar}{2\gamma_0^2mc}\um_0.
 \label{eq:uplus-kn}
\end{align}
这里的一阶自旋量近似要求
\begin{equation}
 \boxed{
 \frac{|\delta k_n|^2}{2}
 \left\|\partial_k^2\widehat u^+(k_0)\right\|
 \ll
 |\delta k_n|\left\|\partial_k\widehat u^+(k_0)\right\|
 \quad\text{且}\quad
 |\delta k_n|\left\|\partial_k\widehat u^+(k_0)\right\|\ll1.}
 \label{eq:uplus-sideband-condition}
\end{equation}
对负能基底同样有
\begin{align}
 \um(k_n)
 &=\um_0+\delta k_n\left.\frac{\partial\um}{\partial k}\right|_{k_0}+O(\delta k_n^2)\\
 &\simeq
 \um_0-\delta k_n\frac{\hbar}{2\gamma_0^2mc}\up_0.
 \label{eq:uminus-kn}
\end{align}
若只取零阶，则
\begin{equation}
 \boxed{\um(k_n)\simeq\um_0}
 \label{eq:uminus-smallrange}
\end{equation}
需要
\begin{equation}
 \boxed{
 |\delta k_n|
 \left\|\partial_k\widehat u^-(k_0)\right\|
 \ll1.}
 \label{eq:uminus-condition}
\end{equation}
这里不能只用 $|n|\hbar\omega/\mathcal E_0\ll1$ 作为所有近似的判据，因为在低速极限下电子静能 $mc^2$ 很大，而纵向反冲仍可能不可忽略。对第 $n$ 阶边带应分别控制：
\begin{equation}
\boxed{
\epsilon_{k,n}\equiv\frac{|\delta k_n|}{k_0}
=\frac{|n|\hbar\omega}{\gamma_0m v_0^2}\ll1,}
\label{eq:small-sideband-k}
\end{equation}
它保证边带动量仍处于中心动量附近；而自旋量基底变化还要求式~\eqref{eq:uplus-sideband-condition} 与 \eqref{eq:uminus-condition}。只有在讨论 $\gamma_n$ 或总能量的相对变化时，才使用
\begin{equation}
\boxed{
\epsilon_{E,n}\equiv\frac{|n|\hbar\omega}{\mathcal E_0}\ll1.}
\label{eq:small-sideband-energy}
\end{equation}

再看相位。由 $z'=z-v_0t$，
\begin{align}
 \ee^{\ii n\kappa z'}\ee^{\ii(k_0z-\Omega_0t)}
 &=\exp\left\{\ii\left[n\kappa(z-v_0t)+k_0z-\Omega_0t\right]\right\}\\
 &=\exp\left\{\ii\left[(k_0+n\kappa)z-(\Omega_0+n\kappa v_0)t\right]\right\}\\
 &=\ee^{\ii(k_nz-\Omega_nt)},
\end{align}
其中
\begin{equation}
 \boxed{
 k_n=k_0+n\frac{\omega}{v_0},
 \qquad
 \Omega_n=\Omega_0+n\omega.}
 \label{eq:kn-omegan}
\end{equation}

这里 $k_n$ 与 $\Omega_n$ 同时按线性关系移动，本质上使用了电子色散在 $k_0$ 附近的一阶展开。把精确正能色散 $E(k)=\sqrt{m^2c^4+\hbar^2c^2k^2}$ 展开：
\begin{align}
 E(k_0+\delta k)
 &=\mathcal E_0
 +\left.\frac{\dd E}{\dd k}\right|_{k_0}\delta k
 +\frac12\left.\frac{\dd^2E}{\dd k^2}\right|_{k_0}\delta k^2+\cdots.
\end{align}
先求一阶导数：
\begin{align}
 \frac{\dd E}{\dd k}
 &=\frac{\hbar^2c^2k}{E}
 =\hbar\frac{pc^2}{E}
 =\hbar v.
\end{align}
再求二阶导数。利用 $\dd E/\dd p=v$ 与
\[
 \frac{\dd v}{\dd p}=\frac{1}{\gamma^3m},
\]
可得
\begin{equation}
 \frac{\dd^2E}{\dd k^2}
 =\hbar^2\frac{\dd v}{\dd p}
 =\frac{\hbar^2}{\gamma^3m}.
 \label{eq:Esecond}
\end{equation}
所以
\begin{equation}
 E(k_0+\delta k)
 =\mathcal E_0+\hbar v_0\delta k
 +\frac{\hbar^2\delta k^2}{2\gamma_0^3m}+O(\delta k^3).
 \label{eq:dispersion-taylor}
\end{equation}
取 $\delta k=\delta k_n=n\omega/v_0$，一阶项正好是
\[
 \hbar v_0\delta k_n=n\hbar\omega.
\]
因此式 \eqref{eq:kn-omegan} 成立的色散线性化条件是二阶项远小于一阶项：
\begin{align}
 \frac{\hbar^2\delta k_n^2}{2\gamma_0^3m}
 &\ll\hbar v_0|\delta k_n|,\\
 \boxed{
 \frac{\hbar|\delta k_n|}{2\gamma_0^3mv_0}
 =\frac{|n|\hbar\omega}{2\gamma_0^3mv_0^2}
 \ll1.}
 \label{eq:linear-dispersion-condition}
\end{align}

用式 \eqref{eq:uplus-kn}，式 \eqref{eq:psi-before-spinor-repack} 先化为
\begin{align}
 \Psi(z,t)
 \simeq\sum_n\left[
 g_n^+(z')\up(k_n)
 -\ii\frac{\hbar}{2\gamma_0^2mc}
 (\partial_{z'}g_n^+)\um_0
 \right]
 \ee^{\ii(k_nz-\Omega_nt)}.
 \label{eq:psi-step1}
\end{align}
在小 $\delta k_n$ 范围内，由式 \eqref{eq:spinor-derivatives} 和 \eqref{eq:uminus-kn}，
\begin{align}
 \left.\frac{\partial\up}{\partial k}\right|_{k_n}
 &=\frac{\hbar}{2\gamma_n^2mc}\um(k_n)\\
 &\simeq\frac{\hbar}{2\gamma_0^2mc}\um_0,
 \label{eq:du-at-kn-approx}
\end{align}
其中还要求
\begin{equation}
 \boxed{
 \left|\frac{\gamma_n-\gamma_0}{\gamma_0}\right|
 \simeq\frac{|n|\hbar\omega}{\mathcal E_0}\ll1.}
 \label{eq:gamma-n-condition}
\end{equation}
因此最终波函数可写成
\begin{equation}
 \boxed{
 \Psi(z,t)\simeq
 \sum_{n=-\infty}^{+\infty}
 \left[
 g_n^+(z')\up(k_n)
 -\ii\,\partial_{z'}g_n^+(z')
 \left.\frac{\partial\up(k)}{\partial k}\right|_{k_n}
 \right]
 \ee^{\ii(k_nz-\Omega_nt)}.}
 \label{eq:final-dirac-wavefunction}
\end{equation}
这就是一阶近似下的完整 Dirac 出射波函数。




\section{非相对论一致性检验：从 Dirac--Pauli 到 Schr\"odinger 包络}

现在从 Schr\"odinger 方程独立推导同一问题。这里不把相对论结果直接 ``令 $c\to\infty$'' 作为答案，而是先从最小耦合哈密顿量出发，再把两条推导逐项比较。


\subsection{从 Dirac 到 Pauli：先消去小分量，再决定是否忽略自旋}

在真正转向标量 Schr\"odinger 方程之前，先从 Dirac 方程本身推出 Pauli 极限，这样可以明确“非相对论”与“忽略自旋”是两个不同步骤。写
\begin{equation}
 \Psi(\bm r,t)=
 \ee^{-\ii mc^2t/\hbar}
 \begin{pmatrix}
 \varphi(\bm r,t)\\
 \chi(\bm r,t)
 \end{pmatrix},
 \qquad
 \bm\pi\equiv-\ii\hbar\bm\nabla-q_{\rm e}\bm A.
 \label{eq:dirac-nr-ansatz}
\end{equation}
代入式~\eqref{eq:dirac-3d}，上、下二分量满足
\begin{align}
 \ii\hbar\partial_t\varphi
 &=q_{\rm e}\phi\,\varphi+c\bm\sigma\cdot\bm\pi\,\chi,
 \label{eq:dirac-large-component}\\
 \ii\hbar\partial_t\chi
 &=(q_{\rm e}\phi-2mc^2)\chi
 +c\bm\sigma\cdot\bm\pi\,\varphi.
 \label{eq:dirac-small-component}
\end{align}
非相对论条件不是把 $\chi$ 直接设为零，而是要求其时间变化与势能都远小于静能隙：
\begin{equation}
 \left\|\ii\hbar\partial_t\chi\right\|,
 \ |q_{\rm e}\phi|\,\|\chi\|
 \ll 2mc^2\|\chi\|.
 \label{eq:small-component-adiabatic}
\end{equation}
于是式~\eqref{eq:dirac-small-component} 的最低阶给出
\begin{equation}
 \boxed{
 \chi\simeq\frac{\bm\sigma\cdot\bm\pi}{2mc}\varphi.}
 \label{eq:small-component-leading}
\end{equation}
将其代回大分量方程：
\begin{equation}
 \ii\hbar\partial_t\varphi
 =\left[
 q_{\rm e}\phi+
 \frac{(\bm\sigma\cdot\bm\pi)^2}{2m}
 \right]\varphi.
 \label{eq:pauli-before-square}
\end{equation}
现在不能把 $(\bm\sigma\cdot\bm\pi)^2$ 简单写成 $\bm\pi^2$，因为不同机械动量分量不对易。利用
\begin{equation}
 \sigma_i\sigma_j=\delta_{ij}+\ii\epsilon_{ijk}\sigma_k,
 \qquad
 [\pi_i,\pi_j]=\ii q_{\rm e}\hbar\epsilon_{ijk}B_k,
\end{equation}
有
\begin{align}
(\bm\sigma\cdot\bm\pi)^2
&=\pi_i\pi_j(\delta_{ij}+\ii\epsilon_{ijk}\sigma_k)\\
&=\bm\pi^2
+\frac{\ii}{2}\epsilon_{ijk}\sigma_k[\pi_i,\pi_j]\\
&=\bm\pi^2
+\frac{\ii}{2}\epsilon_{ijk}\sigma_k
\left(\ii q_{\rm e}\hbar\epsilon_{ij\ell}B_\ell\right)\\
&=\bm\pi^2-q_{\rm e}\hbar\bm\sigma\cdot\bm B,
 \label{eq:sigma-pi-square-pauli}
\end{align}
其中用了 $\epsilon_{ijk}\epsilon_{ij\ell}=2\delta_{k\ell}$。因此 Dirac 方程的最低非相对论极限是 Pauli 方程
\begin{equation}
 \boxed{
 \ii\hbar\partial_t\varphi
 =\left[
 \frac{(-\ii\hbar\bm\nabla-q_{\rm e}\bm A)^2}{2m}
 +q_{\rm e}\phi
 -\frac{q_{\rm e}\hbar}{2m}\bm\sigma\cdot\bm B
 \right]\varphi.}
 \label{eq:pauli-equation}
\end{equation}
对电子 $q_{\rm e}=-e$，Zeeman 项为 $+(e\hbar/2m)\bm\sigma\cdot\bm B$。只有在固定自旋、磁场耦合相对纵向电场相位可忽略时，才进一步退化到下面使用的标量 Schr\"odinger 方程。因此近似层级为
\begin{equation}
 \boxed{
 \text{Dirac}
 \xrightarrow{v/c\ll1}\text{Pauli}
 \xrightarrow{\text{固定自旋且 Zeeman 可忽略}}
 \text{Schr\"odinger}.}
 \label{eq:dirac-pauli-sch-hierarchy}
\end{equation}


\subsection{从三维最小耦合 Schr\"odinger 方程到包络方程}

非相对论最小耦合哈密顿量为
\begin{equation}
H_{\rm NR}
=\frac{1}{2m}\bigl(\widehat{\bm p}-q_{\rm e}\bm A\bigr)^2+q_{\rm e}\phi,
\qquad
\widehat{\bm p}=-\ii\hbar\bm\nabla.
\label{eq:nr-hamiltonian}
\end{equation}
先保留任意 $\phi$ 和三维 $\bm A$。取自由载波
\begin{equation}
\Psi_{\rm NR}(\bm r,t)
=f_{\rm NR}(\bm r,t)
\exp\!\left[\frac{\ii}{\hbar}(\bm p_0\cdot\bm r-\mathcal E_0^{\rm NR}t)\right],
\qquad
\mathcal E_0^{\rm NR}=\frac{p_0^2}{2m},
\qquad
\bm v_0=\frac{\bm p_0}{m}.
\label{eq:nr-vector-carrier}
\end{equation}
定义与相对论部分相同的相对机械动量算符
\(
\bm\Pi=-\ii\hbar\bm\nabla-q_{\rm e}\bm A
\)。由于
\begin{equation}
(\widehat{\bm p}-q_{\rm e}\bm A)
\left[f_{\rm NR}\ee^{\ii\bm p_0\cdot\bm r/\hbar}\right]
=\ee^{\ii\bm p_0\cdot\bm r/\hbar}
(\bm p_0+\bm\Pi)f_{\rm NR},
\end{equation}
所以
\begin{align}
\frac{(\bm p_0+\bm\Pi)^2}{2m}
&=\frac{p_0^2}{2m}
+\frac{\bm p_0}{m}\cdot\bm\Pi
+\frac{\bm\Pi^2}{2m}\notag\\
&=\mathcal E_0^{\rm NR}
+\bm v_0\cdot(-\ii\hbar\bm\nabla-q_{\rm e}\bm A)
+\frac{\bm\Pi^2}{2m}.
\end{align}
自由载波能量严格相消后得到三维非相对论包络方程
\begin{equation}
\boxed{
\ii\hbar(\partial_t+\bm v_0\cdot\bm\nabla)f_{\rm NR}
=\left[
q_{\rm e}(\phi-\bm v_0\cdot\bm A)
+\frac{\bm\Pi^2}{2m}
\right]f_{\rm NR}.}
\label{eq:nr-envelope-vector}
\end{equation}
这与相对论正能包络式~\eqref{eq:rel-envelope-tensor} 的结构完全一致，唯一的自由色散差异是
\begin{equation}
\bm M^{-1}
=\frac{\bm P_\perp}{\gamma_0m}
+\frac{\bm P_\parallel}{\gamma_0^3m}
\quad\longrightarrow\quad
\frac{\bm I}{m}
\qquad(v_0/c\to0),
\label{eq:mass-tensor-low-speed}
\end{equation}
此外 Schr\"odinger 理论没有 Dirac 正负能带与相应的自旋投影修正。这一矢量形式已经直接证明了低速对应。

\subsection{一维显式展开、相对论极限核对与反冲尺度}

再对一维纵向模型逐项展开，以核对各算符系数。

对任意波函数 $\Psi$，
\begin{align}
(\widehat{\bm p}-q_{\rm e}\bm A)^2\Psi
={}&\widehat{\bm p}^{\,2}\Psi
-q_{\rm e}\,\widehat{\bm p}\cdot(\bm A\Psi)
-q_{\rm e}\,\bm A\cdot\widehat{\bm p}\Psi
+q_{\rm e}^2A^2\Psi.
\end{align}
其中
\begin{align}
\widehat{\bm p}\cdot(\bm A\Psi)
&=-\ii\hbar\,\bm\nabla\cdot(\bm A\Psi)\\
&=-\ii\hbar(\bm\nabla\cdot\bm A)\Psi
-\ii\hbar\,\bm A\cdot\bm\nabla\Psi,
\\
\bm A\cdot\widehat{\bm p}\Psi
&=-\ii\hbar\,\bm A\cdot\bm\nabla\Psi.
\end{align}
因此
\begin{equation}
H_{\rm NR}\Psi
=
\left[
-\frac{\hbar^2}{2m}\nabla^2
+\ii\frac{q_{\rm e}\hbar}{m}\bm A\cdot\bm\nabla
+\ii\frac{q_{\rm e}\hbar}{2m}(\bm\nabla\cdot\bm A)
+\frac{q_{\rm e}^2A^2}{2m}+q_{\rm e}\phi
\right]\Psi.
\label{eq:nr-h-expanded}
\end{equation}
在完整三维 Coulomb gauge 中 $\bm\nabla\cdot\bm A=0$。为得到可逐项比较的一维显式式，进一步选择 temporal gauge $\phi=0$，并只保留沿 $z$ 的包络变化与 $A_z$ 分量。此时
\begin{equation}
\ii\hbar\partial_t\Psi
=
-\frac{\hbar^2}{2m}\partial_z^2\Psi
+\ii\frac{q_{\rm e}\hbar}{m}A_z\partial_z\Psi
+\ii\frac{q_{\rm e}\hbar}{2m}(\partial_zA_z)\Psi
+\frac{q_{\rm e}^2A_z^2}{2m}\Psi.
\label{eq:nr-sch-1d}
\end{equation}

把自由载波和慢包络分开：
\begin{equation}
\Psi_{\rm NR}(z,t)
=f_{\rm NR}(z,t)
\exp\!\left[\ii(k_0z-\Omega_0^{\rm NR}t)\right],
\qquad
\hbar\Omega_0^{\rm NR}=\frac{\hbar^2k_0^2}{2m},
\qquad
v_0=\frac{\hbar k_0}{m}.
\label{eq:nr-envelope-ansatz}
\end{equation}
逐项求导：
\begin{align}
\partial_t\Psi_{\rm NR}
&=
\left(\partial_tf_{\rm NR}-\ii\Omega_0^{\rm NR}f_{\rm NR}\right)
\ee^{\ii(k_0z-\Omega_0^{\rm NR}t)},\\
\partial_z\Psi_{\rm NR}
&=
\left(\partial_zf_{\rm NR}+\ii k_0f_{\rm NR}\right)
\ee^{\ii(k_0z-\Omega_0^{\rm NR}t)},\\
\partial_z^2\Psi_{\rm NR}
&=
\left(
\partial_z^2f_{\rm NR}
+2\ii k_0\partial_zf_{\rm NR}
-k_0^2f_{\rm NR}
\right)
\ee^{\ii(k_0z-\Omega_0^{\rm NR}t)}.
\label{eq:nr-derivatives}
\end{align}
代入式 \eqref{eq:nr-sch-1d} 并约去公共相位：
\begin{align}
\ii\hbar\partial_tf_{\rm NR}
+\hbar\Omega_0^{\rm NR}f_{\rm NR}
={}&
-\frac{\hbar^2}{2m}\partial_z^2f_{\rm NR}
-\ii\frac{\hbar^2k_0}{m}\partial_zf_{\rm NR}
+\frac{\hbar^2k_0^2}{2m}f_{\rm NR}\notag\\
&+\ii\frac{q_{\rm e}\hbar}{m}A_z\partial_zf_{\rm NR}
-\frac{q_{\rm e}\hbar k_0}{m}A_zf_{\rm NR}\notag\\
&+\ii\frac{q_{\rm e}\hbar}{2m}(\partial_zA_z)f_{\rm NR}
+\frac{q_{\rm e}^2A_z^2}{2m}f_{\rm NR}.
\end{align}
使用
\[
\hbar\Omega_0^{\rm NR}=\frac{\hbar^2k_0^2}{2m},
\qquad
v_0=\frac{\hbar k_0}{m},
\]
自由载波项严格相消，得到
\begin{align}
\ii\hbar\partial_tf_{\rm NR}
={}&
-\frac{\hbar^2}{2m}\partial_z^2f_{\rm NR}
-\ii\hbar v_0\partial_zf_{\rm NR}
+\ii\frac{q_{\rm e}\hbar}{m}A_z\partial_zf_{\rm NR}\notag\\
&-q_{\rm e}v_0A_zf_{\rm NR}
+\ii\frac{q_{\rm e}\hbar}{2m}(\partial_zA_z)f_{\rm NR}
+\frac{q_{\rm e}^2A_z^2}{2m}f_{\rm NR}.
\end{align}
两边除以 $\ii\hbar$，并把自由传播项移到左侧：
\begin{equation}
\boxed{
\begin{aligned}
\partial_tf_{\rm NR}
+v_0\partial_zf_{\rm NR}
-\frac{\ii\hbar}{2m}\partial_z^2f_{\rm NR}
={}&
\ii\frac{q_{\rm e}v_0}{\hbar}A_zf_{\rm NR}
+\frac{q_{\rm e}A_z}{m}\partial_zf_{\rm NR}\\
&+\frac{q_{\rm e}^2A_z^2}{2\ii\hbar m}f_{\rm NR}
+\frac{q_{\rm e}}{2m}(\partial_zA_z)f_{\rm NR}.
\end{aligned}}
\label{eq:nr-master}
\end{equation}
这就是非相对论包络主方程。若从一开始就在完整三维 Coulomb gauge 中工作，则最后一项由 $\bm\nabla\cdot\bm A=0$ 消失。

现在直接与相对论式 \eqref{eq:rel-master} 比较。令 $v_0/c\ll1$。则
\begin{align}
\gamma_0
&=(1-\left(\frac{v_0}{c}\right)^2)^{-1/2}
=1+\frac12\left(\frac{v_0}{c}\right)^2+\frac38\left(\frac{v_0}{c}\right)^4+O(\left(\frac{v_0}{c}\right)^6),\\
\gamma_0^3
&=1+\frac32\left(\frac{v_0}{c}\right)^2+O(\left(\frac{v_0}{c}\right)^4),\\
\frac{1}{\gamma_0^3}
&=1-\frac32\left(\frac{v_0}{c}\right)^2+O(\left(\frac{v_0}{c}\right)^4).
\label{eq:gamma-low-speed}
\end{align}
因此相对论主方程中所有
\[
\frac{1}{\gamma_0^3m}
\]
都满足
\begin{equation}
\frac{1}{\gamma_0^3m}
=
\frac1m\left[1-\frac32\frac{v_0^2}{c^2}
+O\!\left(\frac{v_0^4}{c^4}\right)\right].
\label{eq:mass-low-speed}
\end{equation}
把式 \eqref{eq:mass-low-speed} 代回式 \eqref{eq:rel-master}。为了保持量纲一致，不把无量纲的 $O(v_0^2/c^2)$ 直接加在有量纲项之后，而是写成相对修正：
\begin{align}
-\frac{\ii\hbar}{2\gamma_0^3m}\partial_z^2f
&=
-\frac{\ii\hbar}{2m}
\left[1-\frac32\frac{v_0^2}{c^2}
+O\!\left(\frac{v_0^4}{c^4}\right)\right]\partial_z^2f,\\
\frac{q_{\rm e}A_z}{\gamma_0^3m}\partial_zf
&=
\frac{q_{\rm e}A_z}{m}
\left[1-\frac32\frac{v_0^2}{c^2}
+O\!\left(\frac{v_0^4}{c^4}\right)\right]\partial_zf,\\
\frac{q_{\rm e}^2A_z^2}{2\ii\hbar\gamma_0^3m}f
&=
\frac{q_{\rm e}^2A_z^2}{2\ii\hbar m}
\left[1-\frac32\frac{v_0^2}{c^2}
+O\!\left(\frac{v_0^4}{c^4}\right)\right]f,\\
\frac{q_{\rm e}}{2\gamma_0^3m}(\partial_zA_z)f
&=
\frac{q_{\rm e}}{2m}
\left[1-\frac32\frac{v_0^2}{c^2}
+O\!\left(\frac{v_0^4}{c^4}\right)\right](\partial_zA_z)f.
\end{align}
主线性耦合
\[
\ii\frac{q_{\rm e}v_0}{\hbar}A_zf
\]
本身不含 $\gamma_0$。于是
\begin{equation}
\boxed{
\text{式 \eqref{eq:rel-master}}
\ \xrightarrow{\,v_0/c\ll1\,}\
\text{式 \eqref{eq:nr-master}},}
\qquad
\text{各相对论系数的相对修正为 }O\!\left(\frac{v_0^2}{c^2}\right).
\label{eq:rel-to-nr}
\end{equation}
这严格说明了：\emph{完整的 Schr\"odinger 包络方程是 Dirac 正能包络方程的低速近似。}

同一结论还可以在自旋量层面看见。由式 \eqref{eq:uplus}，
\begin{align}
\up(k_0)
&=
\begin{pmatrix}
\sqrt{(\gamma_0+1)/(2\gamma_0)}\\
\sqrt{(\gamma_0-1)/(2\gamma_0)}
\end{pmatrix}\\
&=
\begin{pmatrix}
1-\dfrac18\left(\frac{v_0}{c}\right)^2+O(\left(\frac{v_0}{c}\right)^4)\\[0.4em]
\dfrac12\frac{v_0}{c}+O\!\left(\frac{v_0^3}{c^3}\right)
\end{pmatrix}.
\label{eq:spinor-low-speed}
\end{align}
因此 $v_0/c\ll1$ 时，下分量相对于上分量只有 $O(v_0/c)$；电子正能 Dirac 波函数逐渐退化为单分量 Schr\"odinger 波函数。

若进一步进入一阶 PINEM 区域，非相对论方程 \eqref{eq:nr-master} 中各高阶项的控制参数正是相对论条件 \eqref{eq:first-order-allconditions} 在 $\gamma_0\to1$ 时的极限：
\begin{align}
\epsilon_{\rm disp}^{\rm NR}
&\sim\frac{\Delta p_{\rm env}}{2mv_0}\ll1,\\
\epsilon_{\rm grad}^{\rm NR}
&\sim\frac{\Delta p_{\rm env}}{mv_0}\ll1,\\
\epsilon_{A^2}^{\rm NR}
&=\frac{|q_{\rm e}A_z|}{2mv_0}\ll1,
\qquad
\bm\nabla\cdot\bm A=0.
\label{eq:nr-first-order-cond}
\end{align}
于是
\begin{equation}
\boxed{
\partial_tf_{\rm NR}
+v_0\partial_zf_{\rm NR}
\simeq
\ii\frac{q_{\rm e}v_0}{\hbar}A_zf_{\rm NR}.}
\label{eq:nr-first-order}
\end{equation}
令
\[
\zeta=z-v_0t,
\qquad
f_{\rm NR}(z,t)=g_{\rm NR}(\zeta,t),
\]
则
\begin{align}
\partial_tf_{\rm NR}
&=\partial_tg_{\rm NR}-v_0\partial_\zeta g_{\rm NR},\\
\partial_zf_{\rm NR}
&=\partial_\zeta g_{\rm NR},
\end{align}
从而
\begin{equation}
\partial_tg_{\rm NR}
=
\ii\frac{q_{\rm e}v_0}{\hbar}
A_z(\zeta+v_0t,t)g_{\rm NR}.
\end{equation}
固定 $\zeta$ 积分：
\begin{equation}
\boxed{
\frac{g_{\rm NR}(\zeta,t)}
{g_{\rm NR}(\zeta,t_0)}
=
\exp\left[
\ii\frac{q_{\rm e}v_0}{\hbar}
\int_{t_0}^{t}
A_z(\zeta+v_0t',t')\,\dd t'
\right].}
\label{eq:nr-path-phase}
\end{equation}
这与相对论式 \eqref{eq:g-integral} 完全相同。因此在相同的准单色、穿越时间和弱反冲条件下，后续的
\[
\Ftilde\!\left(\frac{\omega}{v_0}\right),
\qquad
\xi_n(\zeta),
\qquad
k_n=k_0+n\frac{\omega}{v_0}
\]
全部与相对论一阶推导相同。非相对论最终波函数只需把 Dirac 自旋量替换成标量载波：
\begin{equation}
\boxed{
\Psi_{\rm NR}^{\rm out}(z,t)
\simeq
\sum_{n=-\infty}^{+\infty}
g_n(\zeta)
\exp\!\left[
\ii\left(
k_nz-\Omega_n^{\rm NR}t
\right)
\right],}
\label{eq:nr-out}
\end{equation}
其中
\begin{equation}
k_n=k_0+n\frac{\omega}{v_0},
\qquad
\Omega_n^{\rm NR}
=\Omega_0^{\rm NR}+n\omega,
\label{eq:nr-sidebands}
\end{equation}
而第二式仍然建立在电子色散的一阶线性化上。对非相对论色散
\[
E_{\rm NR}(p)=\frac{p^2}{2m}
\]
有
\begin{align}
E_{\rm NR}(p_0+\Delta p)-E_{\rm NR}(p_0)
&=v_0\Delta p+\frac{\Delta p^2}{2m}.
\end{align}
取
\[
\Delta p_n=n\hbar\frac{\omega}{v_0},
\]
要求二阶反冲项远小于一阶项：
\begin{equation}
\boxed{
\frac{|n|\hbar\omega}{2mv_0^2}\ll1
\qquad\Longleftrightarrow\qquad
\left|n\frac{\omega}{v_0}\right|\ll k_0.}
\label{eq:nr-recoil}
\end{equation}


当这一反冲条件不再极强时，最先需要恢复的并不是完整强场 Dirac 动力学，而是各边带离开相互作用区后的自由色散相位。对相对论色散，令
\begin{equation}
 \delta k_n=n\kappa=n\frac{\omega}{v_0}.
\end{equation}
由前面的 Taylor 展开，
\begin{align}
E(k_0+\delta k_n)
&=\mathcal E_0
+\hbar v_0\delta k_n
+\frac{\hbar^2\delta k_n^2}{2\gamma_0^3m}
+\frac{1}{6}E_0^{(3)}\delta k_n^3+\cdots\\
&=\mathcal E_0+n\hbar\omega
+\frac{n^2\hbar^2\omega^2}{2\gamma_0^3mv_0^2}
+O(n^3\kappa^3).
 \label{eq:sideband-recoil-dispersion}
\end{align}
因此边带在相互作用结束后自由传播时间 $T$ 会额外获得相对于线性色散的二次相位
\begin{equation}
 \boxed{
 \Phi_n^{(2)}(T)
 =-\frac{n^2\hbar\omega^2}{2\gamma_0^3mv_0^2}T.}
 \label{eq:quadratic-recoil-phase}
\end{equation}
这里已经除去了共同相位 $\mathcal E_0T/\hbar$ 和线性边带相位 $n\omega T$。于是出射谱振幅的一般自由演化是
\begin{equation}
 c_n(T)=c_n(0)
 \exp\!\left[\ii\Phi_n^{(2)}(T)+O(n^3\kappa^3T)\right].
 \label{eq:sideband-free-propagation}
\end{equation}
只测能谱概率时这类相位不改变 $|c_n|^2$；但若随后再经过第二个相干 PINEM 相互作用区，或者把不同边带重新转换为时间/空间密度调制，$n^2$ 相位会产生可观测干涉。一个自然的二次色散时间尺度由 $|\Phi_1^{(2)}|\sim1$ 给出
\begin{equation}
 \boxed{
 T_{\rm rec}
 \equiv\frac{2\gamma_0^3mv_0^2}{\hbar\omega^2}.}
 \label{eq:recoil-time-scale}
\end{equation}
所以“一阶 PINEM 区域”更准确地说是：相互作用区内可忽略色散；若后续漂移距离足够长，即使相互作用本身满足 eikonal 条件，也可能必须恢复自由传播产生的边带二次相位。

需要特别区分两个层次的结论。完整理论上，式 \eqref{eq:nr-master} 的确只是式 \eqref{eq:rel-master} 的 $v_0/c\ll1$ 极限；但是当色散、包络梯度和 $A_z^2$ 等含 $\gamma_0^3m$ 的项全部按照式 \eqref{eq:first-order-allconditions} 舍去以后，两种理论都降到同一个一阶输运方程 \eqref{eq:first-order-pinem} 与 \eqref{eq:nr-first-order}。因此一阶 PINEM 相位与 Bessel 边带公式的形式只显式依赖真实电子速度 $v_0$。因此，即使电子速度已经进入相对论区，只要一阶 eikonal 条件仍成立，两种一阶输运方程仍可具有相同的形式；这种一阶一致性比 ``非相对论只在 $v_0\ll c$ 有效'' 更强，但它不意味着完整 Schr\"odinger 动力学在任意相对论速度下都等价于 Dirac 动力学。





\section{高阶电子有效理论：Foldy--Wouthuysen 展开}

Pauli 方程保留了小分量消元的最低阶。若要判断 Darwin 项、自旋--轨道耦合与相对论动能修正的大小，需要把同一 Dirac 哈密顿量继续作系统展开。写成“偶算符 + 奇算符”的形式
\begin{equation}
H_D=\beta_{\rm D}mc^2+\mathcal E+\mathcal O,
\qquad
\mathcal E=q_{\rm e}\phi,
\qquad
\mathcal O=c\bm\alpha\cdot\bm\pi,
\label{eq:FW-even-odd}
\end{equation}
其中
\begin{equation}
[\beta_{\rm D},\mathcal E]=0,
\qquad
\{\beta_{\rm D},\mathcal O\}=0.
\end{equation}
“偶”表示不混合 Dirac 的大、小分量，“奇”表示会把正、负能块相互耦合。FW 变换的目标就是逐阶消去 $\mathcal O$。

取第一步幺正变换
\begin{equation}
U_1=\ee^{\ii S_1},
\qquad
S_1=-\frac{\ii\beta_{\rm D}\mathcal O}{2mc^2}.
\label{eq:FW-generator}
\end{equation}
因为 $\beta_{\rm D}\mathcal O$ 是反 Hermitian，$S_1$ 是 Hermitian，所以 $U_1$ 确实幺正。对含时变换，新的哈密顿量为
\begin{equation}
H_1=U_1H_DU_1^\dagger+
\ii\hbar(\partial_tU_1)U_1^\dagger.
\label{eq:FW-time-dependent-transform}
\end{equation}
用 Baker--Campbell--Hausdorff 公式逐阶展开：
\begin{align}
U_1H_DU_1^\dagger
={}&H_D+\ii[S_1,H_D]
-\frac12[S_1,[S_1,H_D]]+\cdots,\\
\ii\hbar(\partial_tU_1)U_1^\dagger
={}&-\hbar\partial_tS_1
-\frac{\ii\hbar}{2}[S_1,\partial_tS_1]+\cdots .
\end{align}
最先检查静能项。由 $\{\beta,\mathcal O\}=0$，
\begin{align}
[S_1,\beta_{\rm D} mc^2]
&=-\frac{\ii}{2mc^2}
[\beta_{\rm D}\mathcal O,\beta_{\rm D} mc^2]\\
&=-\frac{\ii}{2}
(\beta_{\rm D}\mathcal O\beta_{\rm D}-\mathcal O)\\
&=-\frac{\ii}{2}(-\mathcal O-\mathcal O)
=\ii\mathcal O,
\end{align}
所以
\begin{equation}
\ii[S_1,\beta mc^2]=-\mathcal O.
\end{equation}
它恰好消去原哈密顿量中的最低阶奇算符。这就是 FW 变换为什么能逐阶“块对角化” Dirac 哈密顿量。

经过继续收集同一相对论阶的项，并把残余奇算符再作一次同类型变换，可以得到“相对于 Pauli 主项的首个 $v^2/c^2$ 修正阶”。为避免把形式上的 $1/m$ 次数与实际的 $v/c$ 阶数混淆，下面保留 $\mathcal O^4$ 动能修正以及同阶的场梯度项：
\begin{align}
H_{\rm FW}
={}&\beta mc^2+\mathcal E
+\frac{\beta_{\rm D}\mathcal O^2}{2mc^2}
-\frac{\beta_{\rm D}\mathcal O^4}{8m^3c^6}
\nonumber\\
&-\frac{1}{8m^2c^4}
\left[
\mathcal O,
[\mathcal O,\mathcal E]+\ii\hbar\dot{\mathcal O}
\right]
+\mathcal R_{\rm FW}^{(\rm higher)}.
\label{eq:FW-general-order2}
\end{align}
其中 $\mathcal R_{\rm FW}^{(\rm higher)}$ 统一表示下一 FW 阶的残余算符；它的具体大小应由 $\|\bm\pi\|/(mc)$、$|q_{\rm e}\phi|/(mc^2)$ 与场变化尺度共同估计，而不能只看一个形式幂次。这里保留 $\dot{\mathcal O}$ 很重要，因为 PINEM 的光场本来就是含时的；它会和 $[\mathcal O,\mathcal E]$ 合并成规范不变的电场。

先计算 $\mathcal O^2$。利用前面已经证明的式~\eqref{eq:sigma-pi-square-pauli}，Dirac 矩阵版本完全相同：
\begin{equation}
\mathcal O^2
=c^2\left(\bm\pi^2-q_{\rm e}\hbar\bm\Sigma\cdot\bm B\right),
\qquad
\bm\Sigma=
\begin{pmatrix}
\bm\sigma&0\\0&\bm\sigma
\end{pmatrix}.
\label{eq:FW-O2}
\end{equation}
再计算电场双对易。由于
\begin{align}
[\mathcal O,\mathcal E]
&=c\alpha_i[\pi_i,q_{\rm e}\phi]
=-\ii q_{\rm e}\hbar c\bm\alpha\cdot\bm\nabla\phi,\\
\ii\hbar\dot{\mathcal O}
&=-\ii q_{\rm e}\hbar c\bm\alpha\cdot\partial_t\bm A,
\end{align}
两者相加为
\begin{equation}
[\mathcal O,\mathcal E]+\ii\hbar\dot{\mathcal O}
=\ii q_{\rm e}\hbar c\bm\alpha\cdot\bm E,
\qquad
\bm E=-\bm\nabla\phi-\partial_t\bm A.
\label{eq:FW-gauge-E}
\end{equation}
现在使用
\begin{equation}
(\bm\alpha\cdot\bm A)(\bm\alpha\cdot\bm B)
=\bm A\cdot\bm B
+\ii\bm\Sigma\cdot(\bm A\times\bm B)
\end{equation}
并保持算符次序：
\begin{align}
[\bm\alpha\cdot\bm\pi,\bm\alpha\cdot\bm E]
={}&\bm\pi\cdot\bm E-\bm E\cdot\bm\pi
+\ii\bm\Sigma\cdot
(\bm\pi\times\bm E-\bm E\times\bm\pi)\\
={}&-\ii\hbar\bm\nabla\cdot\bm E
+\ii\bm\Sigma\cdot
(\bm\pi\times\bm E-\bm E\times\bm\pi).
\end{align}
故
\begin{align}
&\left[
\mathcal O,
[\mathcal O,\mathcal E]+\ii\hbar\dot{\mathcal O}
\right]
\nonumber\\
&\qquad=
q_{\rm e}\hbar^2c^2\bm\nabla\cdot\bm E
+q_{\rm e}\hbar c^2\bm\Sigma\cdot
(\bm E\times\bm\pi-\bm\pi\times\bm E).
\label{eq:FW-double-commutator}
\end{align}
投影到正能二分量块 $\beta_{\rm D}\to+1$、$\bm\Sigma\to\bm\sigma$，得到
\begin{align}
H_{\rm FW}^{(+)}
={}&mc^2+q_{\rm e}\phi
+\frac{\bm\pi^2}{2m}
-\frac{q_{\rm e}\hbar}{2m}\bm\sigma\cdot\bm B
\nonumber\\
&-\frac{\left(\bm\pi^2-q_{\rm e}\hbar\bm\sigma\cdot\bm B\right)^2}
{8m^3c^2}
-\frac{q_{\rm e}\hbar^2}{8m^2c^2}\bm\nabla\cdot\bm E
\nonumber\\
&-\frac{q_{\rm e}\hbar}{8m^2c^2}
\bm\sigma\cdot
(\bm E\times\bm\pi-\bm\pi\times\bm E)
+\cdots .
\label{eq:FW-positive-general}
\end{align}
最后两项分别是 Darwin 修正和 Hermitian 的自旋--轨道耦合。若是静电场 $\bm\nabla\times\bm E=0$，则
\begin{equation}
\bm\pi\times\bm E=-\bm E\times\bm\pi,
\end{equation}
因此自旋--轨道项化为
\begin{equation}
H_{\rm SO}
=-\frac{q_{\rm e}\hbar}{4m^2c^2}
\bm\sigma\cdot(\bm E\times\bm\pi).
\end{equation}
这说明 Pauli 方程并不是 Dirac 的全部非相对论内容，而只是 $1/m$ 的最低阶；当纳米近场梯度足够强、电子速度更高或需要自旋分辨时，应从式~\eqref{eq:FW-positive-general} 判断 Darwin、自旋--轨道及相对论动能项是否可以忽略，而不能简单地把“非相对论”与“标量 Schr\"odinger”画等号。


